\clearpage
\phantomsection
\addcontentsline{toc}{chapter}{附录}

\chapter*{附录}

\section*{附录A~~原型系统典型测试问题示例}

为验证本文所提出的基于大模型的软件运维方法在实际场景中的可用性，原型系统设计了若干典型测试问题。以下给出部分代表性示例。

\noindent 1. 知识问答类问题

\noindent （1）软件运维中常见的故障处理流程包括哪些步骤？

\noindent （2）当服务出现连接超时时，通常应从哪些方面进行排查？

\noindent （3）日志分析在智能运维中的主要作用是什么？

\bigskip
\noindent 2. 状态查询类问题

\noindent （1）当前某服务的错误率是否出现异常？

\noindent （2）过去一小时系统响应时间是否明显升高？

\noindent （3）近一段时间内某模块的资源使用率变化情况如何？

\bigskip
\noindent 3. 故障诊断类问题

\noindent （1）过去一小时支付服务为何响应缓慢？

\noindent （2）某服务频繁出现错误日志的可能原因是什么？

\noindent （3）当系统同时出现高延迟和错误率上升时，应优先排查哪些问题？

\bigskip
上述问题覆盖了知识查询、状态感知和故障诊断等多类场景，可用于验证原型系统在不同任务中的分析能力与交互效果。

\section*{附录B~~根因分析提示模板示例}

为了提高大语言模型在软件运维场景中的输出稳定性与可解释性，本文在原型系统中设计了结构化提示模板。以下为根因分析任务的一类示例模板。

{\ttfamily
你是一个面向软件运维场景的智能分析助手。\\
请结合用户问题、运维知识、日志摘要和监控指标，\\
完成故障分析，并按照以下结构输出：\\
\\
1. 故障现象\\
2. 可能原因\\
3. 判断依据\\
4. 处理建议\\
\\
用户问题：\\
\{用户输入问题\}\\
\\
运维知识：\\
\{检索得到的知识片段\}\\
\\
日志摘要：\\
\{日志分析结果\}\\
\\
监控指标摘要：\\
\{指标分析结果\}
}

\bigskip
通过采用结构化提示方式，系统可以在多源上下文基础上生成更具条理性的分析结果，从而增强根因分析任务中的可解释性与可读性。

\section*{附录C~~系统核心功能说明}

本文设计的原型系统主要包括以下几个核心功能模块。

\noindent （1）用户交互模块：接收用户输入的问题，并展示系统返回结果。

\noindent （2）任务理解模块：识别用户问题属于知识问答、状态查询还是故障诊断。

\noindent （3）知识检索模块：从运维知识库中召回相关文档片段。

\noindent （4）数据处理模块：对日志与监控指标进行提取、筛选和摘要。

\noindent （5）模型分析模块：基于大语言模型完成知识整合、诊断推理与结果生成。

\bigskip
上述模块共同构成了本文所实现的面向软件运维场景的智能分析原型系统。